% ---
% Reemplaza líneas 2026–2231 de Docs/Technical_Report.tex.
% Decisiones que lo respaldan: D-005, D-009, D-013, D-014, D-016, D-017, D-022, D-023, D-024, D-025, D-030, D-031.
% ---

\subsection{Descripción general del proceso}

\begin{figure}[H]
\centering
\begin{tikzpicture}[every node/.style={scale=0.7}, node distance=0.5cm and 0.3cm]

\node[blkGray] (ds) {%
  \textbf{CBIS-DDSM Dataset}\\[2pt]
  {\normalfont\scriptsize
  1{,}566 pacientes $\cdot$ 3{,}568 lesiones $\cdot$ DICOM $+$ CSV\\
  masas y calcificaciones como subproblemas independientes}%
};

\node[blkTeal, below=of ds] (m1) {%
  \textbf{M1 $\cdot$ Lectura de datos}\\[2pt]
  {\normalfont\scriptsize
  \texttt{pydicom} $\cdot$ \texttt{pandas} $\cdot$ DICOM + CSV unificados}%
};

\node[blkTeal, below=of m1] (m2) {%
  \textbf{M2 $\cdot$ Preprocesamiento}\\[2pt]
  {\normalfont\scriptsize
  Alineación y limpieza de máscaras $\cdot$
  Recorte ROI a resolución original $\cdot$ intensidades de 16 bits}%
};

\node[blkTeal, below=of m2] (m3) {%
  \textbf{M3 $\cdot$ Extracción de características}\\[2pt]
  {\normalfont\scriptsize
  PyRadiomics: primer orden, forma, GLCM, GLRLM
  $\rightarrow \mathbf{x} \in \mathbb{R}^{67}$}%
};

\node[blkTeal, below=of m3] (m4) {%
  \textbf{M4 $\cdot$ Selección y escalado}\\[2pt]
  {\normalfont\scriptsize
  F-test con restricción $|r| \le 0.95$ ($k = 12$)
  $+$ cuantiles $\rightarrow \mathbf{x} \in [0,\pi]^{12}$}%
};

\node[font=\tiny\color{colGray}, below=0.25cm of m4] (lm5) {%
  M5 $\cdot$ Transformación (5 condiciones experimentales)%
};
\coordinate (fork) at ($(lm5.south) + (0,-0.2cm)$);

\node[cond, anchor=north] at ($(fork) + (-6.2cm, -0.1cm)$) (c1){%
  \textbf{C1 $\cdot$ Sin}\\ \textbf{transf.}\\[4pt]
  {\normalfont\tiny Baseline}\\ {\normalfont\tiny crudo}%
};
\node[cond, anchor=north] at ($(fork) + (-3.1cm, -0.1cm)$) (c2){%
  \textbf{C2 $\cdot$ PCA}\\ \textbf{lineal}\\[4pt]
  {\normalfont\tiny Control}\\ {\normalfont\tiny nulo}%
};
\node[cond, anchor=north] at ($(fork) + (0cm, -0.1cm)$) (c3){%
  \textbf{C3 $\cdot$ Kernel}\\ \textbf{PCA (RBF)}\\[4pt]
  {\normalfont\tiny Baseline}\\ {\normalfont\tiny no lineal}%
};
\node[condQ, anchor=north] at ($(fork) + (3.2cm, -0.1cm)$) (c4){%
  \textbf{C4 $\cdot$ Qiskit}\\ \textit{ZZ feature map}\\[4pt]
  {\normalfont\tiny ansatz con $\theta$ fijo}\\
  {\normalfont\tiny$\langle Z_i\rangle \in [-1,1]^{12}$ $\cdot$ $K_Q$}%
};
\node[condQ, anchor=north] at ($(fork) + (6.3cm, -0.1cm)$) (c5){%
  \textbf{C5 $\cdot$ Qiskit}\\ \textit{Pauli feature map} $(Y, ZZ)$\\[4pt]
  {\normalfont\tiny ansatz con $\theta$ fijo}\\
  {\normalfont\tiny$\langle Z_i\rangle \in [-1,1]^{12}$ $\cdot$ $K_Q$}%
};

%% Convergencia centrada bajo C3 y a la altura de las cajas cuánticas, que son las más altas
\coordinate (merge) at ($(c3.south |- c5.south) + (0,-0.5cm)$);
\coordinate (mergeL) at (merge -| c1);
\coordinate (mergeR) at (merge -| c5);

%% M6 y M7 en paralelo: ambos consumen la salida de M5
\node[blkAmber, anchor=north, text width=5.6cm, minimum height=1.9cm]
  at ($(merge) + (-2.9cm, -0.6cm)$) (m6){%
  \textbf{M6 $\cdot$ Análisis de separabilidad}\\[2pt]
  {\normalfont\scriptsize
  KTA $\cdot$ \textit{geometric difference} $g_{CQ}$ $\cdot$
  Davies-Bouldin $\cdot$ Fisher $\cdot$ t-SNE}%
};
\node[blkCoral, anchor=north, text width=5.6cm, minimum height=1.9cm]
  at ($(merge) + (2.9cm, -0.6cm)$) (m7){%
  \textbf{M7 $\cdot$ Clasificación (MLP)}\\[2pt]
  {\normalfont\scriptsize
  Mismo MLP en las 5 condiciones $\cdot$
  validación cruzada 5-fold $\cdot$ test oficial}%
};

\node[blkCoral, anchor=north, text width=12cm]
  at ($(m6.south -| merge) + (0,-0.6cm)$) (m8){%
  \textbf{M8 $\cdot$ Evaluación comparativa}\\[2pt]
  {\normalfont\scriptsize
  AUC-ROC $\cdot$ F1 $\cdot$ Accuracy $\cdot$
  correlación separabilidad--clasificación $\cdot$
  masas vs.\ calcificaciones $\cdot$ restricciones del simulador}%
};

\node[blkGray, below=of m8, text width=12cm] (res) {%
  \textbf{Informe de resultados}\\[2pt]
  {\normalfont\scriptsize
  Separabilidad $\cdot$ Clasificación $\cdot$
  Viabilidad práctica del simulador Qiskit Aer}%
};

%% ── Flechas
\foreach \a/\b in {ds/m1, m1/m2, m2/m3, m3/m4, m8/res}
  \draw[arr] (\a.south) -- (\b.north);
\draw[color=colGray, line width=0.4pt] (m4.south) -- (lm5.north);
\draw[color=colGray, line width=0.4pt] (lm5.south) -- (fork);
\draw[color=colGray, line width=0.4pt] (fork -| c1.north) -- (fork -| c5.north);
\foreach \c in {c1, c2, c3, c4, c5}
  \draw[arr] (fork -| \c.north) -- (\c.north);
\foreach \c in {c1, c2, c3, c4, c5}
  \draw[color=colGray, line width=0.4pt] (\c.south) -- (\c.south |- merge);
\draw[color=colGray, line width=0.4pt] (mergeL) -- (mergeR);
\draw[arr] (merge -| m6.north) -- (m6.north);
\draw[arr] (merge -| m7.north) -- (m7.north);
\draw[arr] (m6.south) -- (m6.south |- m8.north);
\draw[arr] (m7.south) -- (m7.south |- m8.north);

\end{tikzpicture}

\caption{%
  Pipeline del sistema híbrido cuántico-clásico para prediagnóstico de cáncer de mama.
  Masas y calcificaciones recorren el pipeline como subproblemas independientes.
  El módulo M4 selecciona 12 características por F-test con una restricción de
  redundancia y las transforma en ángulos en $[0,\pi]$ mediante una transformación por
  cuantiles, ajustada solo sobre el conjunto de entrenamiento; el mismo vector alimenta las
  cinco condiciones. En C4 y C5 los parámetros del \textit{ansatz} se fijan con una semilla
  y no se entrenan, y el kernel de fidelidad $K_Q$ se calcula con el \textit{feature map}
  solo. C4 combina una codificación de primer orden en $Z$ con el acoplamiento $ZZ$, y C5
  una codificación de primer orden en $Y$ con el mismo acoplamiento. La salida de M5 alimenta en paralelo el análisis de
  separabilidad (M6) y la clasificación (M7), que el módulo M8 compara.%
}
\label{fig:pipeline_sistema}
\end{figure}

El flujo es lineal hasta M4 y se ramifica a partir de M5. Las cinco condiciones de M5 reciben el
mismo vector $\mathbf{x}\in[0,\pi]^{12}$, y sus salidas alimentan en paralelo dos módulos
independientes: el análisis de separabilidad (M6), que mide la geometría de cada representación
sin entrenar ningún clasificador, y la clasificación (M7), que entrena el mismo perceptrón
multicapa sobre cada una. El módulo M8 compara ambos resultados. La separabilidad no alimenta a la
clasificación: ambas consumen la salida de M5, y su relación es precisamente lo que M8 evalúa.
